\section {Part 1}

\begin{itemize}
\item {
    \textbf{In your *own* words (please no LLM!), what is the motivation of the paper?}

    HCI researchers often need to set up VR experiments, which, among other things, require human-like hands to increase participants' immersion into the experiment environment. While there are many textured hand models available, there is no standardized, research-oriented open-source solution backed by actual data, which this study \cite{Pohl2022Hands} aims to provide. The discepency between the available hands also may also cause discrepencies in the results of different studies, which might affect their reproducibility. A standardized solution would eliminate these discrepencies and allow researchers to set up more consistent and reproducible experiments.
}

\item {
    \textbf{What are the contributions (to the scientific community) of the paper?}

    The paper provides a set of hand textures for a standard Oculus hand model which were also tested in a user study to ensure that they contribute to participants' immersion. The textures are available for free under MIT license, which makes them a useful tool for HCI researchers who want to set up VR experiments with realistic hands, but either not have a technical know-how, don't have resources to create their own textures or want to avoid discrepencies introduced by different hand textures across studies.
}

\item {
    \textbf{Which research contribution types introduced by Wobbrock and Kientz (2016) for HCI apply to this paper? (Max 2)}

    The study could be clasified as both a Methodological and a Dataset contribution. The former is because the paper literally provides a tool for other researchers to use in their studies, which is a defining characteristic of a Methodological contribution. The latter is because the paper provides a dataset of participants' survey responses related to environment immersion, which is also one of the defining characteristics of a Dataset contribution. The paper also analyzes the data it provides and makes conclusions regarding the affects of hand texture matching on VR immersion, though it is mostly used as justification for the provided hand textures' usefulness and not as a standalone contribution.
}
\end{itemize}
